\section{Why pass to a relative boundary rather than fill every residue?}
The preceding factor-exchange result takes actual bounded affine lines of prime
exponents in $N=p+4u$ and forces their packet to cover an entire cyclic chart.
Its compensating monomials, phase role, integer multiplicities, and inverse are
retained here. Requiring every residue of that chart is sometimes unnecessarily
strong: an available product can fill a proper subgroup after the remaining
coordinates have selected the correct quotient class. The work below supplies
that missing return, not an assumption that a quotient solution lifts.

The source convention is $G(A)=\{\tau\}\sqcup A^\bullet$, with
$e=0_A^\bullet\ne\tau$. A coefficient killed by reduction still has its receiving
support. The scalar $e$ is not the exterior-channel label. The principal source
interface is the actual quotient-to-kernel map BR1 in \cite{SZ}; its application
below is proved directly. No analytic Gamma norm, theta estimate, source-wide
Lean certificate, or unproved arithmetic uniformity is transported into this
finite coefficient calculation.

\section{Original arithmetic and the prescribed ordered return}
Fix
\[
 k\ge3,\quad m=2^k,\quad o=2^{k-2},\quad u=2^{2k-5},\qquad
 p\equiv1\pmod8,\quad p>2u,\quad N=p+4u=\prod_q q^{e_q}. \tag{A1}
\]
The $q$ are distinct actual primes. Structural theorems allow integer $p$;
primality is separately proved in examples. The original middle states at the
fixed seed are in bijection with
\[
 Q\mid N,\quad Q\equiv-1\pmod m,\quad R=N/Q,\quad a=(p+R)/4. \tag{A2}
\]
Indeed $K(u)=\prod_{\ell^f\parallel u}\ell^{\lceil f/2\rceil}=2^{k-2}$ and
$u\mid a^2\iff K(u)\mid a$. Since $m\mid4u$, (A2) implies
$R\equiv-p\pmod m$, giving that availability. Both $Q,R$ are $3\pmod4$, so
at least three; $N<3p$ gives $R,Q<p$. Also $4(a+u)=R+N$ is divisible by $R$.
Reversing these steps proves the correspondence from the original gate and
availability conditions.

With $d_0=\gcd(a,u)$ put
\[
 (h,r,s)=\left(d_0^2/u,u/d_0,a/d_0\right),\qquad
 \lambda=(r+s)/R,\qquad (x,y,z)=(a,phs\lambda,phr\lambda). \tag{A3}
\]
Prime valuations give integrality, $a=hrs$, $u=hr^2$, and $\gcd(r,s)=1$.
Since $\gcd(a,R)=1$, the gate $R\mid a+u=hr(r+s)$ gives $R\mid r+s$.
Multiplying the reciprocal sum by $phrs\lambda$ gives
$(p+R)\lambda=4hrs\lambda$. The ordered inverse is
\[
 \boxed{u=\frac{pa^2}{Ry-pa}.}                              \tag{A4}
\]
The factorization of $N$ is input; the returned factorization of $a$ is distinct
output. The inherited Type II crosswalk is
$(a_{ET},b_{ET},c_{ET},d_{ET},e_{ET},f_{ET})=(r,s,\lambda,h,R,Q)$, with
Elsholtz--Tao denominator order $(x,z,y)$ \cite[\S2, (2.13)--(2.22)]{ET}.

For either $g=3$ or $g=-3$, let $H_g=\langle g\rangle\le(\mathbb Z/m)^\times$.
It has order $o$, contains $p$, and does not contain $-1$.
The two charts consist of residues $1,3\pmod8$ and $1,5\pmod8$, respectively.
Induction from $v_2(3^{2^j}-1)=j+2$ proves these statements and the binary
logarithm lift. All chart logarithms below retain their representatives in $[0,o)$.

\section{The relative norm and its different quotient observations}
Write $o=2^n$, $d=2^v$ with $0\le v<n$, and $L=o/d\ge2$.
In the ordinary coefficient algebra
\[
 A_o=\mathbb F_2[X]/(X^o-1)=\mathbb F_2[T]/(T^o),\quad T=X+1,
\]
define
\[
 \boxed{\mathcal N_d=\sum_{j=0}^{L-1}X^{dj}=T^{o-d}.}        \tag{B1}
\]
The polynomial identity holds before reducing modulo $T^o$: set
$Y=X^d=1+T^d$ and expand the geometric sum by repeated squaring.

\begin{theorem}[Relative boundary comparison and complete inverse]
Let $\pi_d:A_o\to A_d$ send $X$ to its cyclic class of order $d$.
There is an injective $\mathbb F_2[X]$-module map
\[
 \iota_d:A_d\longrightarrow A_o,\qquad X^r\longmapsto X^r\mathcal N_d.
 \tag{B2}
\]
Its image is exactly $\ker(X^d-1)=\ker T^d$. A coefficient vector belongs to
that image precisely when its entries are constant on each congruence class
modulo $d$. Reading the entries at $0,\ldots,d-1$ is its inverse.
Moreover,
\[
 \boxed{\mathcal N_d F=\iota_d\pi_d(F),\qquad
 \pi_d\iota_d=L\,1=0,\qquad \mathcal N_d^2=L\mathcal N_d=0.} \tag{B3}
\]
The first integer identities also hold over $\mathbb Z$, where $L$ is retained.
\end{theorem}
\begin{proof}
The coefficients of $\mathcal N_d F$ at residue $r$ are the sum of all
coefficients of $F$ congruent to $r\pmod d$, proving the first identity.
The shifts $X^r\mathcal N_d$, $0\le r<d$, have disjoint supports, proving
injectivity and the inverse. Multiplication by $X^d-1$ takes successive
coefficient differences along each coset cycle; its kernel is exactly the
constant vectors on those cycles. Each quotient sum of a constant coset is
$L$ times that constant. Multiplying two norm sums counts $L$ decompositions
per subgroup element. These prove all claims, including their integral versions.
\end{proof}

The map (B2) is a module equivalence onto its image, \emph{not} a unital algebra
isomorphism: that image is a square-zero ideal when $L\ge2$. It is the exact
BR1 instance with prequotient $M=\mathbb F_2[T]$, $v_{\rm BR}=T^d$,
and $f=T^{o-d}$. Multiplication by both powers is injective on the prequotient;
the ordinary quotient and the boundary kernel are distinct observations
\cite[BR1--BR2, BR6--BR7]{SZ}.

For $d'\mid d\mid o$, the coefficient repetitions compose:
\[
 \iota_{o,d}\iota_{d,d'}=\iota_{o,d'}.                      \tag{B4}
\]
This follows by the unique expression of a subgroup exponent as $dj+r d'$.
It is not multiplication of two overlapping norm elements in $A_o$.
Over $\mathbb Z$ such a product has the intersection multiplicity
$|H_d\cap H_{d'}|$; over $\mathbb F_2$ it can vanish. The receiving zero keeps
its original support and its integer preimages. No inverse in the global
sparse SplitZero semiring is asserted by the coefficient change $X=T+1$.

\section{A disjoint original-source decomposition, not a rescaled count}
Use a fixed original affine packet as in \cite{Previous}:
\[
 v_q(W)=a_q+\sigma_q t_i\quad(q\in B_i),\quad0\le t_i\le E_i,\qquad
 W(\mathbf t)=A\prod_i G_i^{t_i},                           \tag{C1}
\]
where distinct blocks have disjoint prime indices, each step vector is nonzero,
and the endpoint inequalities place every exponent in $[0,e_q]$.
The $G_i$ can be rational ratios, but the displayed words are actual positive
integer divisors of $N$. Assume $A\notin H_g$, $G_i\in H_g$ and retain
$G_i=g^{a_i}\pmod m$. The inverse reads
$t_i=(v_q(W)-a_q)/\sigma_q$ at a nonzero step and checks all other coordinates.

Include the formal role variable $t_*\in\{0,1\}$ with multiplier $p^{-1}$;
it is never included among the prime factors of $N$. Let $a_*=-\log_g p$.
Together write the step logs as $a_i$ and their capacities as $E_i$.
The fine target is
\[
 \sum_i a_it_i\equiv t_0\pmod o,\qquad
 g^{t_0}=-A^{-1}\pmod m.                                   \tag{C2}
\]
A hit gives $Q=W$ at role zero or $Q=N/W$ at role one, followed by (A3).

For fixed $d=2^v$, define the least strides
\[
 \delta_i=\frac d{\gcd(d,a_i)},\qquad
 t_i=r_i+\delta_i z_i,\quad
 0\le r_i<\delta_i,\quad r_i\le E_i,\quad
 0\le z_i\le F_i(r_i):=\left\lfloor\frac{E_i-r_i}{\delta_i}\right\rfloor.
 \tag{C3}
\]
This is a disjoint decomposition of the entire original parameter box, with
inverse Euclidean division. Identity logs have stride one. It retains every
integer coefficient and every formal role. In tagged generating polynomials,
\[
 \prod_i\sum_{t=0}^{E_i}Z_i^tX^{a_it}
 =\sum_{\mathbf r} Z^{\mathbf r}X^{a\cdot\mathbf r}
   \prod_i\sum_{z=0}^{F_i(r_i)}
          (Z_i^{\delta_i}X^{a_i\delta_i})^z.                \tag{C4}
\]
No inverse-fibre weights are necessary: every original word occurs once.
The original affine prime-coordinate map and its inverse remain attached to
these parameter tags.

\begin{theorem}[Actual relative packet forcing]\label{thm:lift}
Choose a lower digit vector in (C3) satisfying the explicit quotient target
\[
 t_0-a\cdot\mathbf r\equiv0\pmod d.                        \tag{C5}
\]
Set $b_i=a_i\delta_i/d\pmod L$ and, for each nonzero $b_i$,
$w_i=2^{v_2(b_i)}$. If
\[
 \boxed{\sum_{w_i\le2^j}F_i(r_i)w_i\ge2^{j+1}-1,
        \quad0\le j<\log_2L,}                             \tag{C6}
\]
then the original fixed-seed middle branch has a positive integral witness.
All physical exponents, both possible roles, and complete inverse fibres are
retained. The case $d=1$ is the preceding full-chart theorem. The case $d=o$
is excluded: its one-point subgroup would merely restate a found fine target.
\end{theorem}
\begin{proof}
The bounded dyadic selection lemma chooses $c_i\le F_i(r_i)$ with
$\sum_i c_iw_i=L-1$. Necessity is reduction modulo every $2^{j+1}$;
sufficiency follows by selecting a unit coin, pairing the remaining units into
twos, dividing the weights by two, and inducting. Every carry retains its
original item label. A descending-denomination greedy choice follows by exchanging
an available larger coin for a lower submultiset of equal dyadic mass.

In $\mathbb F_2[Y]/(Y^L-1)$, $\operatorname{ord}_{Y+1}(1+Y^{b_i})=w_i$.
Consequently
\[
 \prod_i(1+Y^{b_i})^{c_i}=(Y+1)^{L-1}=\sum_{j=0}^{L-1}Y^j. \tag{C7}
\]
For each $c$, Frobenius gives $(1+U)^c=\sum_{z\preceq c}U^z$ modulo two,
where $z\preceq c$ means binary submask containment. The corresponding
\emph{integer} packet uses each allowed $z_i\preceq c_i$ once. Its coefficient
at every subgroup residue is odd, hence positive. Take the target
$(t_0-a\cdot\mathbf r)/d\pmod L$ and return (C3), (C1), and (A3).
Every endpoint remains within its actual inventory, including $t_*\le1$.

An executable inverse uses binary-factor suffix products. At an odd target
coefficient precisely one of the two next suffix coefficients is odd; choose
that branch. It terminates at zero. Given the state and role, read $W=Q$ or
$R$, invert the affine blocks, divide every parameter by its stride and check
its retained lower digit and high submask. This is the complete fixed-packet
inverse. Different roles can identify only the complementary words $Q,R$.
\end{proof}

Only high steps divisible by $d$ are used in (C7). A mere total $T$-degree
$o-d$ without that restriction is insufficient: in $A_8$,
$(1+X)^5(1+X^3)=T^6+T^7$, not $\mathcal N_2=T^6$.
The relative unit correction has not been discarded.

\paragraph{Uniform fibres, original counts, and computational scope.}
Write $E_i+1=C_i\delta_i+B_i$, $0\le B_i<\delta_i$. On the allowed lower
digits, the complete capacity is
\[
 F_i(r)=C_i-1+\mathbf1_{r<B_i},\qquad
 \min_rF_i(r)=\max(0,C_i-1).                               \tag{C8}
\]
If (C6) holds at these minima, every lower target fibre lifts. The number of
lower target fibres is the actual integer coefficient of (C4) after retaining
only its lower digits and reducing the residue modulo $d$. Each certified
fibre contributes an original word. At a fixed role they are distinct; if both
roles are retained the original state count is at least half their number,
rounded upwards. One cannot sum this bound over different overlapping affine
patterns without recording their further fibres.

For any fixed pattern, the search over $v$ and all its lower digits terminates
and is complete for (C5)--(C6). The computations below use the preceding
canonical maximal singleton and pair lines. Those lines were complete for the
old full-coset test, but different offsets can matter to (C5). No completeness
for \emph{all} relative affine lines is claimed. Arbitrary actual lines remain
valid inputs to Theorem~\ref{thm:lift}. The tests include all old canonical
certificates at $v=0$.

\section{A sharp one-prime lift and a bound on the returned word}
Let $3\le\ell<k$, let $b$ be an actual prime with $v_2(b-1)=\ell$, and write
\[
 N=b^eM,\quad \gcd(b,M)=1,\qquad L=2^{k-\ell},\qquad
 A_\ell(M)=\#\{t\mid M:t\equiv-1\pmod{2^\ell}\}.            \tag{D1}
\]
Then $b$ has order $L$ and generates precisely the principal units
$1\pmod{2^\ell}$. Successive difference-of-squares factorization proves
$v_2(b^{2^j}-1)=\ell+j$, and the two groups have the same size.
For every counted $t$ there is a unique $r_t\in[0,L)$ with
$b^{r_t}\equiv-t^{-1}\pmod{2^k}$.

\begin{theorem}[Complete relative count and uniform capacity]
The entire fixed-seed state count is
\[
 \boxed{C_k(p)=\sum_{\substack{t\mid M\\t\equiv-1\ (2^\ell)}}
 \max\left\{0,1+\left\lfloor\frac{e-r_t}{L}\right\rfloor\right\}.} \tag{D2}
\]
In particular,
\[
 \left\lfloor\frac{e+1}{L}\right\rfloor A_\ell(M)
 \le C_k(p)\le
 \left\lceil\frac{e+1}{L}\right\rceil A_\ell(M).              \tag{D3}
\]
For $e\ge L-1$, $C_k(p)>0$ if and only if $A_\ell(M)>0$.
When $L\mid e+1$, (D3) is equality. The uniform exponent threshold $L-1$
cannot be decreased over the stated integer domain.
\end{theorem}
\begin{proof}
Unique factorization expresses each original cofactor as $Q=b^it$, with
$0\le i\le e$ and $t\mid M$. Its congruence requires the displayed lower
condition and then $i\equiv r_t\pmod L$. This proves the bijection, full fibres
$i=r_t+jL\le e$, and all counts. For sharpness choose distinct primes
\[
 q\equiv-b/3,\qquad r\equiv-b\pmod{2^k},
\]
not equal to $3,b$, and set $M=3qr$, $e=L-2$.
Dirichlet's theorem supplies arbitrarily large actual choices in these reduced
classes \cite[Theorem 18.1]{Dirichlet}. The only divisors of $M$ equal to
$-1\pmod{2^\ell}$ are $r$ and $3q$: both equal $-b\pmod{2^k}$ and require
$i=L-1$. They cannot lift. Nevertheless $N\equiv1\pmod{2^k}$, so
$p=N-4u$ satisfies the original domain once the primes are large enough.
Also $p\equiv1\pmod3$. The derived $p$ is an integer, not asserted prime in
this general sharpness construction.
\end{proof}

For fixed $b,M$, the cofactor count in the family $p_e=b^eM-4u$ has the exact
all-exponent identities
\[
 C(e+L)=C(e)+A_\ell(M),\qquad
 \sum_{e\ge0}C(e)z^e=
 \frac{\sum_t z^{r_t}}{(1-z)(1-z^L)}.                    \tag{D3a}
\]
For each $t$, expand $z^{r_t}/((1-z)(1-z^L))$ to count the same bounded
exponents $i=r_t+jL\le e$. This proves both identities with positive integer
coefficients. The finite prefix for which $p_e\le2u$ is removed explicitly
when interpreting the series as ES states. The primality restriction is a
separate operation; no Euler product or rational series is asserted after
restricting $p_e$ to primes. This family varies $p$, not a fixed-$p$ descent.
The logarithm $r_t$ is computed by binary lifting: at modulus $2^r$ choose
between the previous exponent and its translate by $2^{r-\ell-1}$.
The exact order identity proves uniqueness at each step.

\begin{corollary}[A bounded coarse word gives a bounded original word]
Under $e\ge L-1$ and $A_\ell(M)>0$, an original cofactor can be chosen with
\[
 \boxed{\Omega(Q)\le 2^{k-\ell}-1+2^{\ell-2}.}              \tag{D4}
\]
\end{corollary}
\begin{proof}
Choose a divisor $t\mid M$ equal to $-1\pmod{2^\ell}$ with the fewest prime
occurrences. The quotient of the unit group by $\{1,-1\}$ is cyclic of order
$d=2^{\ell-2}$. If $\Omega(t)>d$, partial sums of its first $d$ prime-log
images give a nonempty subword whose quotient product is one. This is a proper
subword of $t$. Its full residue is $1$ or $-1$. In the latter case retain it;
in the former remove it. Either contradicts minimality. Thus $\Omega(t)\le d$.
The available exponent $r_t\le L-1$ gives $Q=b^{r_t}t$ and (D4).
\end{proof}
The displayed tradeoff can be of order $2^{k/2}$ when an actual prime factor
has $\ell$ near $(k+2)/2$. It does not assert that such a prime or such a
multiplicity occurs for every $p$. This is a cofactor bound, not an unqualified
bound on its complementary residual or a polynomial-time factorization claim.

\section{Strict positive and negative original examples}
Take the certified hard prime
\[
 \boxed{p=8060861761\equiv1\pmod{840},\quad k=6,\quad u=128,}\qquad
 N=3\cdot17^4\cdot53\cdot607.                              \tag{E1}
\]
The full-chart and previous affine-prefix tests fail in both charts. To see
this without a black box, their lowest two weight layers can receive at most
two units: only the primes $3$ and $53$ have odd chart logarithm, each with
capacity one. Pairing can only preserve or consume those two occurrences;
$17$ and $607$ have logarithms divisible by four and eight respectively.
There is no phase item because $p\equiv1\pmod{64}$. The required second prefix
is three. The earlier $3,11$ criterion and all one/two-prime words in either
role fail as well; the verifier independently enumerates their finite domains.

Now $17$ generates $1\pmod{16}$ modulo $64$, of order four.
The only divisors of $M=3\cdot53\cdot607$ equal to $-1\pmod{16}$ are
\[
 t=607,\quad t=159=3\cdot53.
\]
Both are $31\pmod{64}$ and require exponent $r_t=2$.
Thus the complete original branch has exactly two states:
\[
 Q=175423=17^2\cdot607,\qquad Q=45951=17^2\cdot159.          \tag{E2}
\]
For $Q=45951$, the full return is
\[
 (R,a,h,r,s,\lambda)=(175423,2015259296,8,4,62976853,359),
\]
\[
 \boxed{\frac4{8060861761}=\frac1{2015259296}
 +\frac1{1457964212136949677976}+\frac1{92603179910368}.}     \tag{E3}
\]
The other ordered triple is
$(2015226928,5565760054861506408766,353517153390416)$.
Their full factorization and ordered inverse data are serialized.

There is an especially relevant coefficient comparison. In
$\mathbb F_2[C_{16}]$, the four actual choices $17^0,17^1,17^2,17^3$
give $\mathcal N_4=T^{12}$. Its ordinary quotient to $C_4$ is zero.
Further, the two distinct source words $607$ and $159$ have the same residue.
Together they give the \emph{integer} residue counts two at $15,31,47,63$,
which reduce to zero at every occupied coordinate. The original labels and
their two lifted ES states remain. The theorem uses their positive individual
counts, not an inversion of that reduced zero.

The same factor-residue family has an all-multiplicity formula. For primes
$q\equiv53$, $r\equiv31\pmod{64}$, distinct from $3,17$, put
$p_e=3\cdot17^e qr-512$. The only coarse targets are $r,3q$, each requiring
$i\equiv2\pmod4$. Therefore
\[
 \boxed{C_6(p_e)=2\max\left\{0,1+\left\lfloor\frac{e-2}{4}\right\rfloor\right\}.}
 \tag{E4}
\]
This is an exact formula for all integer exponents in the stated domain;
$p_e$ is not asserted prime. At $e\ge2$, $Q=17^2r$ already gives a constructive
positive family. Large binary-encoded $e$ can be counted without expanding $p_e$.

The sharp negative prime is
\[
 \boxed{p=9331009\equiv289\pmod{840},\quad
 N=p+512=3\cdot17^2\cdot229\cdot47.}                        \tag{E5}
\]
There are still two coarse targets, $47$ and $687=3\cdot229$, now both
$47\pmod{64}$. Each requires exponent three while only two are available.
The entire seed-128 branch is empty. This does not disprove ES:
$(p,R,u)=(9331009,7,857)$ gives the middle marking
$(h,r,s,\lambda)=(857,1,2722,389)$ and denominators
$(2332754,8467342993257754,3110706463357)$.

\section{An unbounded separation from fixed short-word cutoffs}
The relative criterion is not just another way to find the three-prime words
in the small census. There is an explicit family on which its advantage over
every fixed occurrence cutoff is unbounded.

\begin{theorem}[Actual separation family]
Fix $k\ge6$, put $L=2^{k-4}$, $b=17$ and $h_*=1+2^{k-1}$.
Choose distinct primes
\[
 q\equiv-h_*/3,\qquad r\equiv-h_*\pmod{2^k},
 \quad N=3\cdot17^Lqr,\quad p=N-2^{2k-3}.                  \tag{E6}
\]
For sufficiently large choices the original domain holds. The complete
fixed-seed branch has exactly two cofactors,
\[
 \boxed{Q_1=17^{L/2}r,\qquad Q_2=3\cdot17^{L/2}q.}         \tag{E7}
\]
The least number of actual prime occurrences in either cofactor/residual role
is $L/2+1$, which tends to infinity with $k$. The previous canonical
singleton/pair full-chart socle tests and the $3,11$ capacity test fail, while
the relative order-$L$ packet succeeds. Auxiliary prime choices can be made
so that the derived integers satisfy $p\equiv1\pmod{840}$; the derived $p$
is not asserted prime without a separate certificate.
\end{theorem}
\begin{proof}
The exact order formula gives $17^{L/2}=h_*$ and $17^L=1\pmod{2^k}$.
The only divisors of $3qr$ equal to $-1\pmod{16}$ are $r,3q$; both have
residue $-h_*$. Hence each requires $i=L/2\pmod L$, with exactly one choice
in $0\le i\le L$. This proves (E7), including completeness and their prime
occurrence lengths. They are complementary, so retaining both roles does not
shorten either word. The high block has capacity $L\ge L-1$ and lifts both
coarse states.

In either old chart, only $3$ and $q$ have odd cyclic logarithm, each with
capacity one. They lie in opposite chart cosets, so cannot form an allowed
pair with each other. Pairing either with another prime has capacity at most
one. The logs of $17$ and $r$ have weights four and $o/2$, so their sums and
differences supply no weight at most two. Consequently the second old prefix
has mass at most two against required three. The phase is zero because
$N\equiv1\pmod{2^k}$. Also $e_3=1,e_{11}=0$, excluding the preceding two-block
capacity rule.

For the hard integer class, also require $q\equiv1\pmod{35}$ and
$r\equiv(1+4u)/(3\cdot17^L)\pmod{35}$. The latter is a unit: $4u$ is an odd
power of two and $1+4u$ is divisible by neither five nor seven. CRT and
Dirichlet supply arbitrarily large primes satisfying these conditions and
(E6). Then $p\equiv1\pmod{35}$, $p\equiv1\pmod{2^k}$, and
$p\equiv1\pmod3$, proving the stated residue. No statement about primality of
the bilinear expression in (E6) is inferred.
\end{proof}

A separately certified prime member of the broader family is
\[
 \boxed{p=16689143913960961\equiv121\pmod{840},\quad k=7,}\qquad
 N=3\cdot17^8\cdot277\cdot2879.                            \tag{E8}
\]
Its cofactors are $240456959=17^4\cdot2879$ and
$69405951=3\cdot17^4\cdot277$. They use five and six occurrences, so even a
four-occurrence search in either role fails. The first returns
\[
 (u,R,a,h,r,s,\lambda)
 =(512,69405951,4172285995841728,8,8,65191968685027,939285).
\]
The two full denominator triples are in the certificate. A complete-order
primality proof uses
\[
 p-1=2^9\cdot3^2\cdot5\cdot13^2\cdot4286125471.
\]
Every displayed prime factor is trial-checked and the required modular order
conditions are verified. No trial division of $p$ through its full square-root
range is claimed for this larger example.

\section{Finite comparison, evidence, and remaining arithmetic obligation}
The finite calculation uses the complete inherited 4519-prime hard universe
through two million and all 37289 valid $k\ge4$ branches. It does not assume
prime occupancy while selecting a packet. The preceding full-affine class
certifies 1443 branches on 1362 primes. Including the nontrivial relative
strides certifies 1723 branches on 1595 primes. Against the stronger union of
that predecessor, short words in both roles, and the $3,11$ criterion, the
actual improvement is
\[
 (3071\text{ branches},2401\text{ primes})
 \longrightarrow(3073\text{ branches},2402\text{ primes}).   \tag{F1}
\]
The two added branches are $(1213129,6)$ and $(1807129,6)$; only $1213129$ is an additional prime for that union. Their high subgroup
has order two. These are explicitly distinguished from the order-four example (E1) and order-eight example (E8). All-divisor occupancy in this seed family remains 3097
branches on 2411 primes and 4583 original states. A further fair baseline is
important: direct words of at most three occurrences already certify 3095
branches on 2409 primes, and their union with the previous tests certifies all
3097 occupied branches on 2411 primes. Thus there is no gain over that stronger
union in this small finite range. The structural separation theorem above and
its prime example (E8) are not inferred from the small census. None of these
counts is new overall ES coverage.

The raw prime-only stride trial was also tested on the 26 occupied branches
missed by the preceding union: it adds none. The affine offsets and exchanges
are retained because they make the strict finite improvement. The canonical
maximal-line search remains a declared class, not an exhaustive solver for
all subfamilies or a proof of universal positivity.

A false prime would require every complete shell to miss, and hence every
fixed-seed branch to evade these actual lift criteria. Failure of one proper
subgroup certificate does not disprove a different certificate. The new
necessary alternative is explicit: each lower target fibre must lack enough
available higher-layer capacity, or its lower target itself must be absent.
The arithmetic correlation forcing some fibre to escape that alternative
across the seeds of one prime is not proved here.

All source maps, integer coefficients, role fibres, missing exponent ranges,
and quotient kernels above are explicit. Finite coefficient calculations and
proof search are not an independent review, a Lean build, or a new analytic
identification. Classical norm and cyclic-cover mechanisms retain attribution
\cite{GH,SZ}; no historical-priority conclusion follows from the scoped search.
